\section{Coalition: who is behind it, and is the opposition manageable?}
\label{sec:coalition}

The seventh question is about the room outside the chamber. A member counts who will
testify for the bill, who will score it, and who, if anyone, will organize against
it.

[DRAFTING INSTRUCTIONS: In the full stage, assemble the coalition from the platform's
stakeholder structure and the bill's standards basis: medical and oncology societies
and patient-safety groups (the safety floor); patient advocates and rural-access
coalitions (the constituent reach); device and robotics makers and standards bodies
(predictable rules of the road, cite \cite{asme2018vv40, isots15066, ul4600}); and
the States, whose human-over-AI laws are preserved. Note the absence of an organized
opponent because the bill creates no entitlement and preempts little. Use
\cite{Kawchak2026PatientPriorityProposedU} on patient priority among proposed bills.
This section carries the coalition mindmap (diagram 13) and the objection-handling
figure (diagram 18).]

\begin{psgfig}
\node[psgact] (c) at (0,0) {Support for\\H. R. 9510};
\node[psgone] (a) at (4.2,1.4) {Clinicians};
\node[psgtwo] (b) at (4.2,0.0) {Patients};
\node[psgthree] (d) at (4.2,-1.4) {Industry};
\draw[psgedge] (c) to[out=40,in=180,looseness=0.6] (a.west);
\draw[psgedge] (c)--(b.west);
\draw[psgedge] (c) to[out=-40,in=180,looseness=0.6] (d.west);
\end{psgfig}
\psgcaption{Figure (placeholder). The coalition radiates into stakeholder groups.
[FULL STAGE: replace with diagram 13; confirm curved arrows carry explicit looseness.]}
